% Standalone problem document, generated from the adjacent README.md.
% Compile with XeLaTeX. Mathematical content is maintained in Markdown.
\documentclass[11pt,a4paper]{article}
\usepackage[fontsize=10.5pt]{scrextend}
\usepackage[margin=22mm,headheight=15pt,headsep=9mm,footskip=11mm]{geometry}
\usepackage{fontspec}
\setmainfont{texgyrepagella}[Extension=.otf,UprightFont=*-regular,BoldFont=*-bold,ItalicFont=*-italic,BoldItalicFont=*-bolditalic]
\setsansfont{texgyreheros}[Extension=.otf,UprightFont=*-regular,BoldFont=*-bold,ItalicFont=*-italic,BoldItalicFont=*-bolditalic]
\setmonofont{texgyrecursor}[Extension=.otf,UprightFont=*-regular,BoldFont=*-bold,ItalicFont=*-italic,BoldItalicFont=*-bolditalic,Scale=MatchLowercase]
\usepackage{amsmath,amssymb}
\usepackage{unicode-math}
\setmathfont{texgyrepagella-math.otf}
\usepackage{xcolor,microtype,enumitem,needspace,fancyhdr}
\usepackage{bookmark}
\definecolor{ink}{HTML}{17344B}
\definecolor{accent}{HTML}{197D80}
\definecolor{muted}{HTML}{596773}
\hypersetup{colorlinks=true,linkcolor=accent,urlcolor=accent,pdftitle={RA-01 - Optimal
pivot count for RPCholesky trace
approximation},pdfauthor={Open Problems in Numerical Linear Algebra}}
\urlstyle{same}
\setlength{\parindent}{0pt}
\setlength{\parskip}{5pt plus 1pt minus 1pt}
\linespread{1.04}
\setlength{\emergencystretch}{3em}
\widowpenalty=10000
\clubpenalty=10000
\displaywidowpenalty=10000
\allowdisplaybreaks[1]
\setlist{nosep,leftmargin=1.5em}
\providecommand{\tightlist}{\setlength{\itemsep}{0pt}\setlength{\parskip}{0pt}}
\providecommand{\pandocbounded}[1]{#1}
\pagestyle{fancy}
\fancyhf{}
\fancyhead[L]{\sffamily\footnotesize\color{muted}OPEN PROBLEMS IN NUMERICAL LINEAR ALGEBRA}
\fancyhead[R]{\sffamily\bfseries\footnotesize\color{accent}RA-01}
\fancyfoot[L]{\parbox[b]{0.9\textwidth}{\sffamily\fontsize{8}{10}\selectfont\color{muted}Editorial ratings; a bounded search cannot certify that a problem remains unsolved.\\Literature check: 2026-09-10}}
\fancyfoot[R]{\sffamily\footnotesize\color{muted}\thepage}
\renewcommand{\headrulewidth}{0.3pt}
\renewcommand{\headrule}{\hbox to\headwidth{\color{accent}\leaders\hrule height \headrulewidth\hfill}}
\makeatletter
\renewcommand{\section}{\Needspace{5\baselineskip}\@startsection{section}{1}{0pt}{12pt plus 2pt minus 2pt}{4pt}{\sffamily\bfseries\large\color{ink}}}
\renewcommand{\subsection}{\Needspace{4\baselineskip}\@startsection{subsection}{2}{0pt}{9pt plus 1pt minus 1pt}{3pt}{\sffamily\bfseries\normalsize\color{ink}}}
\makeatother
\setcounter{secnumdepth}{0}
\begin{document}
{\sffamily\small\color{muted}Randomized and low-rank approximation\par}
\vspace{3pt}
{\raggedright\hyphenpenalty=10000\exhyphenpenalty=10000\sffamily\bfseries\fontsize{21}{24}\selectfont\color{ink}Optimal
pivot count for RPCholesky trace approximation\par}
\vspace{7pt}
{\sffamily\small\textbf{Difficulty:} challenging\quad\textbf{Importance:} interesting
to the community\par}
{\sffamily\small\bfseries\color{accent}Status: Open\par}
\vspace{4pt}
{\color{accent}\hrule height 0.8pt}
\vspace{6pt}
\textbf{Rating rationale:} Challenging because a uniform near-optimal
pivot count must exploit adaptive residual structure; community impact
is efficient kernel and PSD matrix approximation.\\
\textbf{Topic:} randomized low-rank approximation; kernel matrices

\subsection{Context and notation}\label{context-and-notation}

For a Hermitian positive-semidefinite \(A\in\mathbb C^{n\times n}\),
define the exact-arithmetic RPCholesky residuals by \(R_0=A\).
Conditional on \(R_t\ne0\), select \(j\) with probability
\((R_t)_{jj}/\operatorname{tr}(R_t)\) and set

\[
R_{t+1}=R_t-\frac{R_t(:,j)R_t(j,:)}{(R_t)_{jj}}.
\]

If \(R_t=0\), keep all subsequent residuals zero. Eigenvalues are
ordered \(\lambda_1(A)\geq\cdots\geq\lambda_n(A)\geq0\), and
\(\tau_r(A)=\sum_{j>r}\lambda_j(A)\). This is the pivot rule in Chen,
Epperly, Tropp, and Webber,
\href{https://doi.org/10.1002/cpa.22234}{\emph{Randomly pivoted
Cholesky: Practical approximation of a kernel matrix with few entry
evaluations}}, Algorithm 1; their Lemma 5.5 gives the current comparison
\(\mathbb E\operatorname{tr}(R_r)\leq2^r\tau_r(A)\).

\subsection{Problem statement}\label{problem-statement}

Does there exist a universal constant \(C\geq1\) such that, for every
\(n\geq1\), every Hermitian positive-semidefinite
\(A\in\mathbb C^{n\times n}\), every integer \(1\leq r\leq n\), and
every \(0<\varepsilon<1\), RPCholesky satisfies

\[
\mathbb E\operatorname{tr}(R_k)\leq(1+\varepsilon)\tau_r(A),
\qquad k=\min\{n,\lceil Cr/\varepsilon\rceil\}?
\]

The expectation is over the algorithm's adaptive pivots. The constant
must be independent of dimension, spectrum, rank, and tolerance. The
target concerns this fixed pivoting rule.

\subsection{References}\label{references}

Epperly, \href{https://arxiv.org/html/2608.20633v1}{\emph{A new analysis
of the randomly pivoted Cholesky algorithm}}, §1.2, conjecture
immediately after Eq. (1.5), and Corollary 1.3. Earlier motivation
appears in Epperly,
\href{https://tropp.caltech.edu/dissertations/Epp25-Making-Most.pdf}{\emph{Make
the Most of What You Have}}, Caltech dissertation (2025), §11.1, p.~181.

\subsection{Status check}\label{status-check}

The August 2026 paper proves the larger bound
\(k\geq r/\varepsilon+2r\sqrt{\log r}+r\log(1/\varepsilon)+2.3r\) and
explicitly conjectures the displayed improvement. Searches for the
paper's title, \textit{RPCholesky optimal r epsilon}, and
\textit{RPCholesky conjecture} found no subsequent resolution. The
arXiv record listed only v1, posted August 21, 2026. The weaker
dissertation Conjecture 11.1 is therefore excluded as resolved by this
later preprint.

\subsection{Audit --- 2026-09-10}\label{audit-2026-09-10}

Rechecked the \href{https://arxiv.org/html/2608.20633v1}{August 2026
conjecture after (1.5)} and its current arXiv record, still v1.
Optimal-pivot and RPCholesky follow-up searches found no resolution. The
proved extra rank-dependent term does not establish the uniform
\(O(r/\varepsilon)\) target.
\end{document}
